\documentclass[exercice,a4]{thedocument}%
\startdatakeys
    \source{laclassecooperativedestef.com}
    \BO{2026}
    \cycle{3}
    \competencesDuSocle{RE4,CO2}
    \objectifsApprentissage{C3C1ab,C3C1aj,C3C1ad,C3C2ac,C3C1af}
    \motsCles{parallelisme, cercle, quadrilatere}
\enddatakeys
\begin{document}%
\noindent\textbf{Programme de construction 19}\par\smallskip\noindent
\begin{itemize}
    \item Tracer une droite $(\mathspoint{d})$, puis une droite
    $(\mathspoint{e})$ parallèle à $(\mathspoint{d})$ et située au-dessus,
    à $3\text{ cm}$.
    \item Placer le point $\mathspoint{A}$ sur la droite $(\mathspoint{e})$.
    \item Tracer le cercle $\mathcal{C}$ de centre $\mathspoint{A}$ et de
    rayon $3{,}5\text{ cm}$.
    \item Le cercle $\mathcal{C}$ coupe la droite $(\mathspoint{e})$ en
    $\mathspoint{B}$ et $\mathspoint{C}$, et la droite
    $(\mathspoint{d})$ en $\mathspoint{D}$ et $\mathspoint{E}$.
    \item Tracer la figure $\mathspoint{BCDE}$.
    \item Quel est le nom de cette figure ?
\end{itemize}
\end{document}%
